\section{Forschungsansatz: qualitative Einzelfallstudie}
\label{sec:04-01_case-study-approach}

Die vorliegende Untersuchung ist als qualitative Einzelfallstudie konzipiert. Einzelfallstudien eignen sich insbesondere zur Analyse komplexer, zeitlich und organisatorisch eingebetteter Phänomene, bei denen Untersuchungsgegenstand und Kontext nicht unabhängig voneinander betrachtet werden können \cite{404:Yin2018}. Der untersuchte Fall ist der neu entwickelte Wiki-Auftritt einer Produktorganisation innerhalb eines großen Automobilunternehmens, der im Rahmen eines Transformationsprogramms als zentrale Wissens- und Kollaborationsplattform eingeführt wurde.

Ziel der Einzelfallstudie ist es, zu analysieren, wie sich der neue Wiki-Auftritt im Vergleich zum bisherigen Wiki hinsichtlich Nutzbarkeit, Verständlichkeit, Rollenpassung und Akzeptanz bewährt. Im Sinne der Fallstudiensystematik nach \citeauthor{404:Yin2018} steht dabei kein Anspruch statistischer Generalisierbarkeit im Vordergrund. Stattdessen zielt die Untersuchung auf ein vertieftes Verständnis des konkreten Falls sowie der darin wirksamen Gestaltungs- und Nutzungsmechanismen ab \cite{404:Yin2018}. Die Fallstudie dient damit der analytischen Prüfung und Präzisierung theoretischer Annahmen zu UX-Strategien, Informationsarchitektur und Technologieakzeptanz im organisationalen Kontext.

Methodisch handelt es sich um eine qualitative Einzelfallstudie mit eingebetteten Analyseebenen (*embedded case study*) nach Yin \cite{404:Yin2018}. Dabei wird ein übergeordneter Fall durch mehrere analytisch unterscheidbare, jedoch inhaltlich miteinander verbundene Untersuchungseinheiten vertieft analysiert. Abbildung~\ref{fig:g-11_embedded-case-study_yin} veranschaulicht die Struktur des untersuchten Falls sowie dessen Einbettung in drei Analyseebenen.

\image{H}{\defaultImageSize}{g-11_embedded-case-study_yin}{
    Eingebettete Einzelfallstudie des Wiki-Auftritts einer Produktorganisation mit Makro-, Meso- und Mikroanalyseebenen (eigene Darstellung \acrshort{iAa} \citeauthor{404:Yin2018})}

Auf der Makroebene wird der organisationale Kontext betrachtet, insbesondere die Rolle des Wikis innerhalb des Transformationsprogramms sowie seine Funktion als strategische Plattform der Produktorganisation. Die Mesoebene fokussiert die Systemebene und umfasst den Vergleich zwischen dem bisherigen und dem neu eingeführten Wiki, etwa in Bezug auf Informationsarchitektur, Funktionalitäten und strukturelle Gestaltung. Die Mikroebene richtet den Blick auf die konkrete Nutzung des Wikis durch unterschiedliche Nutzerrollen und Personas und stützt sich auf empirische Daten aus verschiedenen Erhebungsformen.

Diese eingebetteten Analyseebenen umfassen im Einzelnen:
\begin{itemize}
    \item unterschiedliche Nutzerrollen und Personas innerhalb der Produktorganisation,
    \item den systematischen Vergleich zwischen altem und neuem Wiki,
    \item sowie verschiedene qualitative und quantitative Erhebungsformen (Usability-Tests, leitfadengestützte Interviews, standardisierte Feedbackumfrage und Hintergrundabfragen).
\end{itemize}

Durch diese mehrstufige Struktur wird es möglich, individuelle Nutzungserfahrungen, systemische Gestaltungsentscheidungen und organisationale Rahmenbedingungen nicht isoliert, sondern in ihrer Wechselwirkung zu analysieren. Die Einbettung mehrerer Analyseebenen dient dabei nicht der Fragmentierung des Falls, sondern der analytischen Verdichtung und Kontextualisierung des Untersuchungsgegenstands im Sinne der von Yin beschriebenen Fallstudienlogik \cite{404:Yin2018}.

Ergänzend zur qualitativen Perspektive wurden standardisierte \Gls{usability}-Messgrößen erhoben, darunter Bearbeitungszeiten, Klickanzahl, Erfolgsraten sowie subjektive Bewertungen auf Ratingskalen. Diese quantitativen Daten sind im Sinne eines Mixed-Methods-Ansatzes unterstützend eingebunden \cite{200:ScribbrInductiveDeductive2023}. Sie dienen der Strukturierung, Veranschaulichung und Plausibilisierung der qualitativ gewonnenen Erkenntnisse, nicht jedoch der statistischen Hypothesenprüfung. Die Auswertung folgt damit einem überwiegend qualitativen, fallbezogenen Erkenntnisinteresse, das punktuell durch deskriptive Kennzahlen ergänzt wird.
